\documentclass[11pt,a4paper]{article}
\usepackage[T1]{fontenc}
\usepackage[utf8]{inputenc}
\usepackage{lmodern}
\usepackage[margin=2.2cm]{geometry}
\usepackage{booktabs}
\usepackage{longtable}
\usepackage{tabularx}
\usepackage{array}
\usepackage{enumitem}
\usepackage{xcolor}
\usepackage[hidelinks]{hyperref}
\usepackage{caption}
\setlist{nosep, leftmargin=*}
\captionsetup{font=small, labelfont=bf}
\newcommand{\code}[1]{\texttt{\small #1}}

\title{DS588 Project -- Phase 1\\[2pt] \large Data Collection and Preparation: TEFAS Funds and US-Listed ETFs}
\author{Burak Kara}
\date{October 2026}

\begin{document}
\maketitle

\section{Overview}
We compare Turkish investment funds traded on TEFAS with US-listed ETFs.
The Phase~1 deliverable is a reproducible Python pipeline. It selects a universe of
\textbf{50 TEFAS investment funds} and \textbf{50 US-listed ETFs}, downloads their
\textbf{daily prices from 3~January~2022 to 5~October~2026} and converts all prices to
\textbf{US dollars}. For the TEFAS funds it also records the trading and settlement terms
needed to simulate realistic orders.
The data were collected on 5~October~2026; the TEFAS trading terms were collected on
12~October~2026.

\section{Universe selection}
\paragraph{Turkish funds (TEFAS).}
TEFAS (Türkiye Elektronik Fon Alım Satım Platformu, operated by Takasbank) lists three
fund types: 1{,}071 investment funds (\code{YAT}), 311 pension funds (\code{EMK}) and
37 exchange-traded funds (\code{BYF}).
We restrict the universe to \code{YAT} funds, for two reasons:
\begin{itemize}
  \item Pension funds can only be bought inside a private pension contract.
  \item The \code{BYF}s are exchange-traded funds listed on Borsa Istanbul, not TEFAS mutual
        funds, so they are kept out to separate ETFs cleanly from other investment funds.
\end{itemize}
The TEFAS price API also returns about 2{,}000 \code{YAT} funds, including roughly
1{,}000 private (\emph{özel}) funds that are not traded on TEFAS. Many of these have a
single investor yet hold billions of TRY, so they would dominate a size ranking. We
therefore keep only the funds flagged as traded on TEFAS (\code{tefasDurum = true}) and
rank them by total net assets (\emph{portföy büyüklüğü}) on 5~October~2026. The top~50
hold 2.91~trillion TRY. Fund sizes range from 14.2 to 192.1~billion TRY.

\paragraph{US ETFs.}
The stockanalysis.com ETF screener lists 5{,}701 US-listed ETFs with assets under management
(AUM). We cross-checked this list against the official Nasdaq ETF screener (5{,}254
symbols, of which 5{,}196 match). We rank by AUM and keep the top~50, which hold
\$8.77~trillion. AUM ranges from \$50~billion (XLF) to \$1.05~trillion (VOO).

\paragraph{Universe composition.}
Table~\ref{tab:composition} summarises the two universes. Selecting by size gives a liquid,
investable universe in both markets, but the two universes differ in kind:
\begin{itemize}
  \item \textbf{TEFAS side:} dominated by FX-denominated \emph{serbest} (hedge-type) funds
        and money-market funds.
  \item \textbf{US side:} mostly equity index ETFs.
\end{itemize}
Full lists are given in Appendix~\ref{app:lists}.

\begin{table}[h]
\centering\small
\caption{Composition of the selected universe (size as of 5 Oct 2026).}
\label{tab:composition}
\begin{tabular}{@{}lr@{\hspace{2.5em}}lr@{}}
\toprule
\multicolumn{2}{c}{\textbf{TEFAS (50 funds)}} & \multicolumn{2}{c}{\textbf{US ETFs (50)}}\\
\cmidrule(r){1-2}\cmidrule(l){3-4}
%%TEFAS_COMPOSITION%%
\bottomrule
\end{tabular}
\end{table}

\section{Instrument identification}
\begin{itemize}
  \item \textbf{TEFAS funds} are open-ended mutual funds, not ETFs. They do not trade on
        an exchange; units are created and redeemed at the daily net asset value (NAV)
        through TEFAS. Each fund is identified by its three-letter TEFAS code and its ISIN.
        The ISIN is stored in \code{tefas\_top50\_trading\_terms.csv}.
  \item \textbf{US instruments} are exchange-traded products identified by their ticker.
        Most are open-ended ETFs. The exceptions are QQQ, which is a unit investment trust,
        and GLD, IAU and IBIT, which are grantor trusts holding physical gold or bitcoin.
        All of them trade intraday on US exchanges, so we treat them as ETFs for
        simulation purposes.
\end{itemize}

\section{Data sources}
\begin{table}[h]
\centering\small
\caption{Data sources and access method.}
\begin{tabularx}{\linewidth}{@{}lXl@{}}
\toprule
\textbf{Data} & \textbf{Source / access} & \textbf{Tool}\\
\midrule
TEFAS fund list, TEFAS status & \code{tefas.gov.tr/api/funds/fonGetiriBazliBilgiGetir} & \code{tefas-crawler}, \code{requests}\\
TEFAS daily NAV, size, investors & \code{tefas.gov.tr/api/funds/fonGnlBlgSiraliGetir} & \code{pytefas} 0.4.1\\
TEFAS trading terms (ISIN, cut-off, settlement) & TEFAS fund detail pages (\code{/tr/fon-detayli-analiz/<code>}) & browser session\\
US ETF list and AUM & stockanalysis.com ETF screener & \code{requests}\\
US ETF list (cross-check) & Nasdaq ETF screener API & \code{requests}\\
US ETF daily prices & Yahoo Finance & \code{yfinance} 1.7.0\\
USD/TRY exchange rate & Yahoo Finance, ticker \code{TRY=X} & \code{yfinance}\\
\bottomrule
\end{tabularx}
\end{table}

All sources are free and publicly accessible, and none requires authentication. TEFAS
limits each request to about one month of data and serves at most five years of history,
so the pipeline requests the data in 28-day chunks.

\section{Pipeline}
The pipeline consists of three scripts:
\begin{itemize}
  \item \code{fetch\_tefas.py} builds the full TEFAS fund list with sizes.
  \item \code{fetch\_us\_etfs.py} builds the full US ETF list with AUM.
  \item \code{fetch\_prices.py} is the main script. Its parameters are \code{TOP\_N=50},
        \code{START=2022-01-01} and \code{ADJUSTED=True}.
\end{itemize}
Running \code{python fetch\_prices.py} reproduces all price files in about 10~minutes. It
performs the following steps:
\begin{enumerate}
  \item Rank the US ETFs by AUM, select the top~$N$ and download their daily closes from
        Yahoo Finance.
  \item Take the TEFAS-traded \code{YAT} funds, rank them by their latest portfolio size
        and select the top~$N$.
  \item Download the TEFAS NAV history in 28-day chunks and keep only the selected funds.
  \item Download USD/TRY, align it to TEFAS dates and divide the TRY NAVs by the rate.
  \item Save each market to a separate wide-format CSV file (date $\times$ instrument).
\end{enumerate}

\section{Prices, currency and adjustments}
\begin{itemize}
  \item \textbf{Frequency and period.} The series are daily, from 3~Jan~2022 to 5~Oct~2026.
        There are 1{,}195 TEFAS dates and 1{,}193 US trading days.
  \item \textbf{TEFAS prices.} The price is the official per-unit NAV published by TEFAS
        once per business day, in TRY. FX-denominated funds (``Döviz'' or ``Döviz-Avro''
        in the name) also publish their NAV in TRY. Their USD series start at about
        \$1.00 for USD funds and about \$1.10 for EUR funds, and grow at roughly 4--5\% per
        year, which confirms that the conversion is correct.
  \item \textbf{Currency conversion.} We use $P^{\text{USD}}_t = P^{\text{TRY}}_t / \text{USDTRY}_t$,
        where USDTRY is the Yahoo Finance daily close. On a TEFAS date with no FX quote we
        use the last available rate (forward fill). The Yahoo rate is not the official
        rate that fund managers use, so the conversion adds a small amount of day-to-day
        noise. It does not bias long-horizon returns.
  \item \textbf{Distributions, TEFAS.} TEFAS funds are accumulating: income is reinvested
        and shows up in the NAV, so the NAV series is already a total-return series. The
        exception is YNK, a dividend-paying (\emph{kâr payı ödeyen}) fund. Its payouts are
        not in the NAV series, so its total return is understated (see
        Section~\ref{sec:limits}).
  \item \textbf{Distributions, US.} We use the Yahoo \emph{adjusted} close
        (\code{auto\_adjust=True}), which is adjusted for splits and for dividends
        (dividends are reinvested on the ex-date). Returns computed from this series are
        total returns. Dividends must therefore \textbf{not} be added again in the
        simulation, as this would double-count them.
  \item \textbf{Fees.} Management fees and ETF expense ratios are already deducted in the
        NAV and the market price. None of the 50 TEFAS funds charges an entry or exit
        commission.
\end{itemize}

\section{Missing observations and cleaning}
\begin{itemize}
  \item \textbf{Late inception.} 28 of the 50 TEFAS funds and one ETF (IBIT, listed
        11~Jan~2024) start after 3~Jan~2022. Their earlier cells are left empty
        (\code{NaN}). This accounts for all 11{,}270 missing TEFAS cells and all 508
        missing US cells. We do not backfill or interpolate. No series has gaps after its
        first observation.
  \item \textbf{Calendars.} The Turkish and US holiday calendars differ: 44 dates exist
        only in the TEFAS file and 42 only in the US file. The two files are kept separate
        and are not forward-filled onto a common calendar. Aligning them is a modelling
        choice for Phase~2.
  \item \textbf{Invalid prices.} Rows with a non-positive price are dropped.
  \item \textbf{Flagged anomaly.} PTE starts at 0.01~TRY at its launch (29~Apr~2025) and
        jumps to about 1~TRY on 12~May~2025, a factor of about 100 in one day. This is a
        unit-price artefact, not a return. The observations before 12~May~2025 should be
        dropped before computing returns. This step is not yet applied in the saved file.
\end{itemize}

\section{Trading, pricing and settlement characteristics}
\paragraph{TEFAS funds.}
TEFAS funds use \emph{forward pricing}: a buy or sell order is not executed at the last
published NAV. It is executed at the next NAV calculated after the order. Orders entered
after the fund's cut-off time are treated as orders for the next business day.

Cash and fund units then settle after a fund-specific number of business days, the
\emph{valör}. Table~\ref{tab:terms} summarises these terms for our 50 funds, as published
on each fund's TEFAS page. The full table is in
\code{data/tefas\_top50\_trading\_terms.csv}.

\begin{table}[h]
\centering\small
\caption{Order cut-off and settlement lag (business days) of the 50 TEFAS funds.}
\label{tab:terms}
\begin{tabular}{@{}cccrl@{}}
\toprule
Cut-off & Buy valör & Sell valör & Funds & Typical type\\
\midrule
%%TERMS%%
\bottomrule
\end{tabular}
\end{table}

Three further restrictions affect who can trade these funds and when:
\begin{itemize}
  \item \textbf{Closed funds.} Two funds, PRY and PSE, are currently closed to both
        subscriptions and redemptions on TEFAS.
  \item \textbf{Qualified investors only.} 33 of the 50 funds are \emph{serbest} funds,
        which are legally restricted to qualified investors. A retail investor could not
        buy them, which must be taken into account when the results are interpreted.
  \item \textbf{Minimum amount.} All 50 funds have a minimum order of one unit.
\end{itemize}

\paragraph{US ETFs.}
ETFs trade continuously on the exchange from 9:30 to 16:00~ET at market prices that can
deviate slightly from NAV. Trades pay the bid--ask spread, which is small for these very
large ETFs. Settlement is T+1 from 28~May~2024 onward; before that date it was T+2.

\paragraph{Proposed simulation convention.}
To avoid assuming that every fund can be traded immediately at an observed price:
\begin{itemize}
  \item A decision taken with information up to the close of day~$t$ is executed for
        TEFAS funds at the NAV of day~$t{+}1$, the next price after the order.
  \item The cash or units from a TEFAS order become available only after the fund's
        buy or sell valör.
  \item ETF orders are executed at the close of $t{+}1$, minus an assumed half-spread.
  \item Funds closed to trading are excluded from the dates on which they were closed.
\end{itemize}

\section{Limitations and reproducible alternatives}\label{sec:limits}
\begin{itemize}
  \item \textbf{Survivorship and look-ahead bias.} The universe is selected by size on the
        \emph{last} date, so funds that shrank or closed during 2022--2026 are excluded.
        A point-in-time alternative is to rank by \code{portfolio\_size} (TEFAS) or AUM on
        the first date of each rebalancing period. The pipeline already downloads
        \code{portfolio\_size} for every date, so this requires no new data source.
  \item \textbf{Unofficial endpoints.} The TEFAS JSON API and the stockanalysis.com
        screener are undocumented and may change. For the US list, the Nasdaq screener
        (symbols only) is an alternative. For TEFAS, the monthly bulletins on the KAP
        (Public Disclosure Platform) fund pages are an alternative.
  \item \textbf{History limit.} TEFAS serves only five years of history. Our start date
        of 1~Jan~2022 must be moved forward once it falls outside that window. To preserve
        the series, the saved CSV files should be archived.
  \item \textbf{Trading terms are a snapshot.} The fund detail pages are protected by a
        JavaScript bot challenge, so the trading terms were read through a browser session
        on 12~Oct~2026 and stored as a snapshot. The terms may have differed earlier in the
        sample. Each fund's prospectus (\emph{izahname}) on KAP is an authoritative,
        reproducible alternative source.
  \item \textbf{Dividend-paying fund.} YNK's distributions are not captured. They can be
        added from the fund's KAP dividend announcements, or the fund can be excluded.
\end{itemize}

\section{Data dictionary}
All files are in \code{data/}. The \code{date} column uses the format YYYY-MM-DD.
\begin{table}[h]
\centering\small
\begin{tabularx}{\linewidth}{@{}lX@{}}
\toprule
\textbf{File} & \textbf{Content and columns}\\
\midrule
\code{us\_top50\_prices\_usd.csv} & Rows are trading dates and columns are the 50 ETF tickers. Values are adjusted daily closes in USD.\\
\code{tefas\_top50\_prices\_usd.csv} & Rows are TEFAS dates and columns are the 50 fund codes. Values are the NAV per unit in USD (TRY NAV / USDTRY).\\
\code{us\_top50.csv} & \code{symbol}, \code{name}, \code{assetClass}, \code{aum\_usd} (AUM in USD, 5~Oct~2026).\\
\code{tefas\_top50.csv} & \code{fund\_code}, \code{fund\_name}, \code{portfolio\_size} (TRY), \code{investor\_count}, \code{date} (date the size was measured).\\
\code{tefas\_top50\_trading\_terms.csv} & \code{fund\_code}, \code{isin}, \code{platform\_status}, \code{risk\_value}, \code{order\_start}, \code{order\_cutoff} (local time), \code{buy\_settlement\_days}, \code{sell\_settlement\_days}, \code{min\_buy\_units}, \code{min\_sell\_units}, \code{entry\_fee}, \code{exit\_fee}, \code{category}, \code{scraped\_on}.\\
\code{tefas\_funds.csv} & All 1{,}419 TEFAS funds (\code{YAT}, \code{EMK} and \code{BYF}) with type, size (\code{portBuyukluk}, TRY), number of investors, market share and returns.\\
\code{us\_etfs.csv} & All 5{,}701 US ETFs with \code{aum\_usd}, \code{assetClass}, \code{price}, \code{volume}, \code{expenseRatio}, \code{inceptionDate}, \code{issuer} and \code{in\_nasdaq\_list}.\\
\bottomrule
\end{tabularx}
\end{table}

\appendix
\section{Selected instruments}\label{app:lists}
\small
\begin{longtable}{@{}rlllrccc@{}}
\caption{TEFAS universe. Size is total net assets in billion TRY (5 Oct 2026). Type includes the fund's denomination (TRY, USD, EUR or gold). ``First'' is the first observation, shown when it is later than 3 Jan 2022.}\\
\toprule
\# & Code & Manager & Type & Size & Valör (B/S) & Cut-off & First\\
\midrule\endfirsthead
\toprule
\# & Code & Manager & Type & Size & Valör (B/S) & Cut-off & First\\
\midrule\endhead
%%TEFAS_TABLE%%
\bottomrule
\end{longtable}

\begin{longtable}{@{}rllp{8.4cm}r@{}}
\caption{US ETF universe. AUM in billion USD (5 Oct 2026).}\\
\toprule
\# & Ticker & Class & Name & AUM\\
\midrule\endfirsthead
\toprule
\# & Ticker & Class & Name & AUM\\
\midrule\endhead
%%US_TABLE%%
\bottomrule
\end{longtable}

\end{document}
